\subsection*{A sharper positive Fourier complement at the certified window}
The previous bound loses constants by bounding each place on the full
space. On the omitted Fourier tail one can retain the limiting
archimedean commutator, the physical incidence of prime shifts, and the
decay of the negative pole vector. These are bounds for the same Weil
operator and its original Fourier projection.

For $L=2\log\lambda$ put
\begin{align}
 R_L&=\frac{e^{-L/2}}{1-e^{-2L}},&
 C_L&=\max\{1,1/4+R_L\},\notag\\
 E_L(t)&=\frac1{15t}+\frac7{2t^2}
       +\frac\pi{2Lt}+\frac{2+2R_L}{Lt^2},\qquad t>0.
 \label{eq:v115-arch-errors}
\end{align}
Define the weighted prime degree
\begin{equation}\label{eq:v115-prime-degree}
 M_\lambda=\mathop{\rm ess\,sup}_{0<x<L}
 \left\{\sum_{\substack{1<m<\lambda^2\\\log m\le x}}
                    \frac{\Lambda(m)}{\sqrt m}
       +\sum_{\substack{1<m<\lambda^2\\\log m\le L-x}}
                    \frac{\Lambda(m)}{\sqrt m}\right\}.
\end{equation}
The strict upper cut omits only a possible full-length translation,
which is zero on $L^2([0,L])$.

\begin{proposition}[A practical positive omitted complement]
\label{prop:v115-sharp-tail}
For $N\ge1$ let $t_*=2\pi(N+1)/L$, and set
\begin{align}
 \Gamma_{\lambda,N}^{\rm even}
   &=\log\frac{N+1}{L}-M_\lambda
                -\frac{2C_L}{t_*}-E_L(t_*),\notag\\
 \Gamma_{\lambda,N}^{\rm all}
   &=\Gamma_{\lambda,N}^{\rm even}-\frac\pi2
                   -\frac{4\sinh^2(L/4)L}{\pi^2N}.
 \label{eq:v115-sharp-tail}
\end{align}
For every form-domain vector $f$ whose Fourier coefficients vanish on
$|n|\le N$,
\[
 QW_\lambda(f,f)\ge\Gamma_{\lambda,N}^{\rm all}\|f\|_2^2.
\]
For an inversion-even $f$, the stronger lower bound with
$\Gamma_{\lambda,N}^{\rm even}$ holds. In particular, at the actual
window $\lambda=3$,
\begin{equation}\label{eq:v115-lambda3-tail}
 \Gamma_{3,256}^{\rm all}>0.0148,\qquad
 \Gamma_{3,256}^{\rm even}>1.5867,\qquad
 \Gamma_{3,512}^{\rm all}>0.7085.
\end{equation}
These inequalities control the entire omitted complement, not merely a
finite further block.
\end{proposition}
\begin{proof}
Write $a_k=2k+1/2$ and use the exact Fourier-grid series for $S_n$
in the proof of Proposition~\ref{prop:v114-weil-core}.
For $t>0$, the function $x\mapsto t/(x^2+t^2)$ decreases on $x>0$.
Comparing its values on $1/2+2\mathbb N$ with its integral gives
\[
 \frac\pi4-\frac1{4t}
 \le\sum_{k\ge0}\frac{t}{a_k^2+t^2}
 \le\frac\pi4+\frac1t.
\]
The exponentially weighted part of $S_n$ is nonnegative and at most
$R_L/|t_n|$. Oddness therefore proves
\begin{equation}\label{eq:v115-S-sharp}
 \left|S_n-\frac\pi4\operatorname{sgn}n\right|
       \le\frac{C_L}{|t_n|}\qquad(n\ne0).
\end{equation}
On the tail, the archimedean off-diagonal multiplier is consequently
$b_n^{\rm arch}=\operatorname{sgn}(n)/4+e_n$ with
$\|e\|_\infty\le C_L/(\pi t_*)$.
Since the compressed discrete Hilbert transform has norm at most $\pi$,
its commutator has norm at most $\pi/2+2C_L/t_*$.
For even coefficients the limiting commutator is instead positive:
under the normalized identification with positive indices its matrix is
\[
 K_{nm}=\frac1{2(n+m)},\qquad n,m>N.
\]
Positivity follows from
$1/(n+m)=\int_0^1x^{n+m-1}\,dx$. Thus only $2C_L/t_*$ is lost in
the even sector. This uses the unshifted Fourier basis of
\eqref{eq:v114-weil-decomposition}; translation to centered physical
coordinates conjugates by $(-1)^n$ and preserves both parity and the
positive quadratic form.

Next retain the complete archimedean diagonal. Its exact series, already
used in the independent matrix certificate, gives for
$z=1/4+it/2$ and $t=|t_n|>0$,
\begin{equation}\label{eq:v115-arch-exact}
 -A_n=\Re\psi_0(z)-\log\pi+\frac{\Re\psi_1(z)}{2L}
 -\frac2L\sum_{k\ge0}e^{-a_kL}
                   \frac{a_k^2-t^2}{(a_k^2+t^2)^2}.
\end{equation}
Here $\psi_0$ and $\psi_1$ are digamma and trigamma.
Binet's exact integral formula~\cite[eq.~5.9.15]{NISTDigamma} at
$w=z+1=5/4+it/2$ has remainder of modulus at most
\[
 2\int_0^\infty\frac{u\,du}
 {|u^2+w^2|(e^{2\pi u}-1)}
 \le\frac8{5t}\int_0^\infty\frac{u\,du}{e^{2\pi u}-1}
 =\frac1{15t},
\]
since $|u^2+w^2|\ge5t/4$. The recurrence
$\psi_0(z)=\psi_0(z+1)-1/z$, together with
$\Re\log w\ge\log(t/2)$,
$\Re(1/(2w))\le5/(2t^2)$, and $\Re(1/z)\le1/t^2$, yields
\[
 \Re\psi_0(z)\ge\log(t/2)-\frac1{15t}-\frac7{2t^2}.
\]
The convergent trigamma series and a decreasing-integrand comparison give
\[
 |\psi_1(z)|\le\sum_{k\ge0}\frac1{(k+1/4)^2+t^2/4}
                 \le\frac4{t^2}+\frac\pi t.
\]
The last term in \eqref{eq:v115-arch-exact} has modulus at most
$2R_L/(Lt^2)$. Hence
\begin{equation}\label{eq:v115-diagonal-sharp}
 -A_n\ge\log\frac{|n|}{L}-E_L(|t_n|).
\end{equation}
Its right side increases with $|n|$, so $|n|=N+1$ supplies the uniform
tail bound. No logarithmic constant or finite-$L$ diagonal term was
discarded from the exact identity.

For the prime part, set $w_m=\Lambda(m)/\sqrt m$ and $a_m=\log m$.
The elementary inequality $2\Re(z\overline w)\le|z|^2+|w|^2$ gives
\[
 -2\sum_m w_m\Re\int_0^{L-a_m}
      f(x+a_m)\overline{f(x)}\,dx
 \ge-\int_0^L d_\lambda(x)|f(x)|^2dx
 \ge-M_\lambda\|f\|_2^2,
\]
where $d_\lambda$ is the function inside braces in
\eqref{eq:v115-prime-degree}. This applies equally to complex vectors.

In centered physical coordinates the negative pole term is
$-2|\ip f{\sinh(x/2)}|^2$. Direct Fourier integration gives
\[
 |\widehat{\sinh(x/2)}(n)|^2
 =\frac{4\sinh^2(L/4)t_n^2}{L(t_n^2+1/4)^2}.
\]
Consequently
\[
 2\|Q_N\sinh(x/2)\|_2^2
 \le\frac{4\sinh^2(L/4)L}{\pi^2}
                         \sum_{n>N}n^{-2}
 \le\frac{4\sinh^2(L/4)L}{\pi^2N}.
\]
It vanishes identically on the even sector; the positive pole term can
be retained or omitted in a lower bound. Combining these estimates
proves \eqref{eq:v115-sharp-tail} first on finite Fourier sums and then
on the tail form domain by Proposition~\ref{prop:v114-weil-core}.

For the numerical constants at $\lambda=3$, symmetry reduces the prime
degree to $0<x<\log3$. Its breakpoints, expressed in $e^x$, are
$9/8,9/7,9/5,2,9/4$. With
$P=w_2+w_3+w_4+w_5+w_7+w_8$, the six interval values are
\[
 P,\quad P-w_8,\quad P-w_8-w_7,\quad
 w_2+w_3+w_4,\quad2w_2+w_3+w_4,\quad2w_2+w_3.
\]
All are at most $P$, using $w_5>w_2$ for the fifth value, and the first
occurs on an interval of positive length. Thus $M_3=P$ exactly.
Also $\sinh^2(L/4)=1/3$, $R_L=27/80$ and $C_L=1$.
Outward interval evaluation of the displayed elementary expressions
gives \eqref{eq:v115-lambda3-tail}; the reproducibility bundle contains
the 256-bit scalar certificate. These are not new finite-matrix
eigenvalue computations.
\end{proof}

At $\lambda=3$ this replaces the earlier sufficient tail cutoff by a
practical Fourier complement. It does not prove positivity of the
entire Weil form: the finite head and its signed correction
$F-B^*T^{-1}B$ must still be controlled on the same operator. The
verified-solve identity \eqref{eq:v114-schur-residual} remains the next
test, with the improved even or full tail constant as appropriate.
